\lesson{3}{Mar 18 2022 Fri (16:53:32)}{Factor sum/difference of two cubes}{Unit 3}
\label{les_3_unit_3:factor_sum_difference_of_two_cubes}

\begin{formula}[Sum of Perfect Cubes]
  \label{frm:sum_of_perfect_cubes}

  \begin{equation}
    \label{eqt:sum_of_perfect_cubes}
    (a^{3} + b^{3}) = (a + b)(a^{2} - ab + b^{2})
  .\end{equation}
\end{formula}

\begin{formula}[Difference of Perfect Cubes]
  \label{frm:difference_of_perfect_cubes}

  \begin{equation}
    \label{eqt:difference_of_perfect_cubes}
    (a^{3} - b^{3}) = (a - b)(a^{2} + ab + b^{2})
  .\end{equation}
\end{formula}

\begin{info}
  It's always helpful if you memorize the table for all of the cubes:

  \begin{tabular}{cc}
    $1^{3} = 1$ & $7^{3} = 343$ \\
    $2^{2} = 8$ & $8^{3} = 512$ \\
    $3^{8} = 27$ & $9^{3} = 729$ \\
    $4^{8} = 64$ & $10^{3} = 1000$ \\
    $5^{8} = 125$ & \\
    $6^{8} = 216$ & \\
  \end{tabular}

  It will help you recognize when you have a sum/difference of two cubes.
\end{info}

\begin{example}
  \label{exm:difference_of_cubes}

  $x^{3} - 27$. We can clearly see that this is a difference of two cubes. Now,
  we need to get the $a$ and $b$ values. To do this, take the cube root of each
  of the variables/numbers, and assign them respectfully. So,
  $a = \sqrt[3]{x^{3}} = x \qquad b = \sqrt[3]{27} = 3$

  Then, plug-in play: $(x - 3)(a^{2} + 3x + 9)$
\end{example}

\begin{example}
  \label{exm:sum_of_cubes}

  $8y^{3} + 1$. We can clearly see that this is a sum of two cubes. Now, let's
  assign the $a$ and $b$ variables. So,
  $a = \sqrt[3]{8y^{3}} = 2y \qquad b = \sqrt[3]{1} = 1$

  Then, plug-in: $(2y + 1)(4y^{2} - 2y + 1)$
\end{example}

\newpage
